\documentclass[a4paper,12pt]{article}
\usepackage[utf8]{inputenc}
\usepackage[T1]{fontenc}
\usepackage[italian]{babel}
\usepackage{pgfplots}
\usepackage{amsmath}
\usepackage{float}
\usepackage{booktabs}
\usepackage{geometry}
\geometry{margin=2cm}

\pgfplotsset{compat=1.18}

\title{Report 4.6 -- Combinazioni di $k$ e TIMES}

\begin{document}

\maketitle

\newpage

\section{Grafici dei risultati per $k$ e $TIMES$ = 10}

I grafici mostrano rispettivamente l'andamento di:
\begin{itemize}
    \item Avg fQ
    \item fQ Min Found
    \item Profit Found
    \item Weight Found
    \item Profit GAP
    \item Weight GAP
    \item n of fQ Min Found
\end{itemize}
al variare di $k \in \{1000, 2000, 3000, 4000, 5000\}$, confrontando
i risultati usando QALS e NON utilizzando QALS.

I risultati riportati sono stati ottenuti utilizzando:

\begin{itemize}
    \item $\lambda = \frac{\sum_{i=1}^{n} p_i}{650 \cdot C}$
    \item l'istanza \texttt{ksp\_1.txt}
\end{itemize}

\begin{table}[H]
\centering
\begin{tabular}{lc}
\toprule
Profit ottimo & 198 \\
Weight ottimo & 101 \\
TIMES & 10 \\
\bottomrule
\end{tabular}
\caption{Soluzione di riferimento del problema di knapsack.}
\end{table}

\section{Tabella riassuntiva}

\begin{table}[H]
\centering
\resizebox{\textwidth}{!}{%
\begin{tabular}{c|rrrrrrr|rrrrrrr}
\toprule
\multirow{2}{*}{$k$} & \multicolumn{7}{c|}{QALS} & \multicolumn{7}{c}{NO QALS} \\
\cmidrule(lr){2-8} \cmidrule(lr){9-15}
 & Avg fQ & fQ Min & Profit & Weight & Profit GAP & Weight GAP & n fQ Min
 & Avg fQ & fQ Min & Profit & Weight & Profit GAP & Weight GAP & n fQ Min \\
\midrule
1000 & -122.17 & -202.98 & 137 & 110 & 61 & -9 & 1 & -330.85 & -330.85 & 302 & 177 & -104 & -76 & 10 \\
2000 & -113.81 & -155.43 & 93 & 76 & 105 & 25 & 1 & -330.85 & -330.85 & 302 & 177 & -104 & -76 & 10 \\
3000 & -145.48 & -195.84 & 131 & 85 & 67 & 16 & 2 & -330.85 & -330.85 & 302 & 177 & -104 & -76 & 10 \\
4000 & -146.27 & -221.06 & 158 & 78 & 40 & 23 & 1 & -330.85 & -330.85 & 302 & 177 & -104 & -76 & 10 \\
5000 & -117.29 & -211.00 & 156 & 59 & 42 & 42 & 1 & -330.85 & -330.85 & 302 & 177 & -104 & -76 & 10 \\
\bottomrule
\end{tabular}%
}
\caption{Risultati numerici per $k \in \{1000,2000,3000,4000,5000\}$ con \texttt{TIMES = 10}.}
\end{table}

\newpage

\section{Grafici}

\pgfplotsset{
    every axis/.append style={
        width=0.92\textwidth,
        height=0.42\textwidth,
        grid=both,
        xlabel={$k$},
        xtick={1000,2000,3000,4000,5000},
        legend style={at={(0.5,-0.22)}, anchor=north, legend columns=2},
        tick label style={font=\small},
        label style={font=\small},
        title style={font=\bfseries}
    }
}

\begin{figure}[H]
\centering
\begin{tikzpicture}
\begin{axis}[title={Avg fQ}, ylabel={Avg fQ}]
\addplot[blue, mark=*] coordinates {(1000,-122.17) (2000,-113.81) (3000,-145.48) (4000,-146.27) (5000,-117.29)};
\addlegendentry{QALS}
\addplot[red, dashed, mark=square*] coordinates {(1000,-330.85) (2000,-330.85) (3000,-330.85) (4000,-330.85) (5000,-330.85)};
\addlegendentry{NO QALS}
\end{axis}
\end{tikzpicture}
\caption{Andamento di Avg fQ al variare di $k$.}
\end{figure}

\begin{figure}[H]
\centering
\begin{tikzpicture}
\begin{axis}[title={fQ Min Found}, ylabel={fQ Min Found}]
\addplot[blue, mark=*] coordinates {(1000,-202.98) (2000,-155.43) (3000,-195.84) (4000,-221.06) (5000,-211.00)};
\addlegendentry{QALS}
\addplot[red, dashed, mark=square*] coordinates {(1000,-330.85) (2000,-330.85) (3000,-330.85) (4000,-330.85) (5000,-330.85)};
\addlegendentry{NO QALS}
\end{axis}
\end{tikzpicture}
\caption{Andamento di fQ Min Found al variare di $k$.}
\end{figure}

\begin{figure}[H]
\centering
\begin{tikzpicture}
\begin{axis}[title={Profit Found}, ylabel={Profit Found}]
\addplot[blue, mark=*] coordinates {(1000,137) (2000,93) (3000,131) (4000,158) (5000,156)};
\addlegendentry{QALS}
\addplot[red, dashed, mark=square*] coordinates {(1000,302) (2000,302) (3000,302) (4000,302) (5000,302)};
\addlegendentry{NO QALS}
\addplot[black, dotted, thick] coordinates {(1000,198) (5000,198)};
\addlegendentry{Profit ottimo}
\end{axis}
\end{tikzpicture}
\caption{Andamento del profitto trovato al variare di $k$.}
\end{figure}

\begin{figure}[H]
\centering
\begin{tikzpicture}
\begin{axis}[title={Weight Found}, ylabel={Weight Found}]
\addplot[blue, mark=*] coordinates {(1000,110) (2000,76) (3000,85) (4000,78) (5000,59)};
\addlegendentry{QALS}
\addplot[red, dashed, mark=square*] coordinates {(1000,177) (2000,177) (3000,177) (4000,177) (5000,177)};
\addlegendentry{NO QALS}
\addplot[black, dotted, thick] coordinates {(1000,101) (5000,101)};
\addlegendentry{Weight ottimo}
\end{axis}
\end{tikzpicture}
\caption{Andamento del peso trovato al variare di $k$.}
\end{figure}

\begin{figure}[H]
\centering
\begin{tikzpicture}
\begin{axis}[title={Profit GAP}, ylabel={Profit GAP}]
\addplot[blue, mark=*] coordinates {(1000,61) (2000,105) (3000,67) (4000,40) (5000,42)};
\addlegendentry{QALS}
\addplot[red, dashed, mark=square*] coordinates {(1000,-104) (2000,-104) (3000,-104) (4000,-104) (5000,-104)};
\addlegendentry{NO QALS}
\addplot[black, dotted, thick] coordinates {(1000,0) (5000,0)};
\addlegendentry{GAP nullo}
\end{axis}
\end{tikzpicture}
\caption{Andamento del Profit GAP al variare di $k$.}
\end{figure}

\begin{figure}[H]
\centering
\begin{tikzpicture}
\begin{axis}[title={Weight GAP}, ylabel={Weight GAP}]
\addplot[blue, mark=*] coordinates {(1000,-9) (2000,25) (3000,16) (4000,23) (5000,42)};
\addlegendentry{QALS}
\addplot[red, dashed, mark=square*] coordinates {(1000,-76) (2000,-76) (3000,-76) (4000,-76) (5000,-76)};
\addlegendentry{NO QALS}
\addplot[black, dotted, thick] coordinates {(1000,0) (5000,0)};
\addlegendentry{GAP nullo}
\end{axis}
\end{tikzpicture}
\caption{Andamento del Weight GAP al variare di $k$.}
\end{figure}

\begin{figure}[H]
\centering
\begin{tikzpicture}
\begin{axis}[title={n of fQ Min Found}, ylabel={n of fQ Min Found}]
\addplot[blue, mark=*] coordinates {(1000,1) (2000,1) (3000,2) (4000,1) (5000,1)};
\addlegendentry{QALS}
\addplot[red, dashed, mark=square*] coordinates {(1000,10) (2000,10) (3000,10) (4000,10) (5000,10)};
\addlegendentry{NO QALS}
\end{axis}
\end{tikzpicture}
\caption{Numero di volte in cui è stato trovato il minimo fQ al variare di $k$.}
\end{figure}

\end{document}